\ifdefined\zoneNsMassInterior\else\InputIfFileExists{96_optimization_num}{}{}\fi

\section{Results on Both Maps}
\label{sec:results}

\input{61_black_holes}
\input{62_neutron_stars}

\subsection{The Winners and Their Zones}
\label{sec:results-zones}

The winner of each pair is scored in its zones beside the data ceiling in \cref{tab:zones}, and \cref{tab:tuning} gives the verdict of each candidate's tuning. On the black holes the winner of the mass, a \resultsBhMassFamily{} surrogate, keeps $\zoneBhMassInterior$ figures in the interior against a ceiling of \resultsBhMassCeiling{}, and the edge \zoneBhMassMargin{} with $\zoneBhMassEdge$ figures; on the charge it keeps $\zoneBhChargeInterior$ against \resultsBhChargeCeiling{}, and the edge \zoneBhChargeMargin{} with $\zoneBhChargeEdge$ figures. On the neutron stars the winner of the mass keeps $\zoneNsMassInterior$ figures against \resultsNsMassCeiling{}, and the edge \zoneNsMassMargin{} with $\zoneNsMassEdge$ figures; on the charge it keeps $\zoneNsChargeInterior$ against \resultsNsChargeCeiling{}, and the edge \zoneNsChargeMargin{} with $\zoneNsChargeEdge$ figures, so the extrapolation did not hold there. Every other candidate ran its search inside the time budget, verdict searched, and none is cut: too slow; the pointwise interpolant is scored on fixed neighbour counts on the neutron-star mass (verdict \tunNsMassRbfPointwiseVerdict{}), on the neutron-star charge (\tunNsChargeRbfPointwiseVerdict{}) and on the black-hole charge (\tunBhChargeRbfPointwiseVerdict{}), and the residual network on the neutron-star charge (\tunNsChargeResnetPointwiseVerdict{}).

\InputIfFileExists{96_optimization_tab_zones}{}{}
\InputIfFileExists{99_precision_tab_tuning}{}{}
